% ============================================================
% CONCLUSION GÉNÉRALE
% ============================================================
\chapter*{Conclusion générale}
\phantomsection\addcontentsline{toc}{chapter}{Conclusion générale}
\markboth{Conclusion générale}{Conclusion générale}

Le Mini-Projet~17 demandait de détecter, sur un flux de virements, des
fraudes qui n'existent qu'à l'échelle du réseau. La plateforme livrée le fait
de bout en bout : un producteur réaliste publie virements et fraudes dans
Kafka ; un job Spark Structured Streaming maintient un graphe glissant de
10~minutes et l'analyse à deux vitesses avec GraphFrames ; chaque alerte part
dans Kafka avec une explication en français et rejoint, avec son sous-graphe,
un lakehouse Delta partitionné par niveau de risque et par date ; un tableau
de bord Streamlit/PyVis permet de l'investiguer.

Les résultats sont mesurés, pas illustrés : contre la vérité terrain publiée
par le générateur, 99,0\,\% des alertes significatives sont fondées,
92,9\,\% des fraudes sont détectées, un carrousel est signalé en
3,9~secondes en médiane, et un arrêt brutal de Spark ne laisse aucun doublon
dans le lakehouse. L'étude d'ablation montre que chaque filtre de la
détection a un effet propre, et l'historique des runs montre d'où viennent
ces chiffres.

Le résultat le plus instructif est sans doute le chemin parcouru en
performance : de 97 à 3,4~secondes par micro-batch, non par un réglage
miracle mais par une suite de mesures, dont une a conduit à annuler une
optimisation. La leçon dépasse ce projet : en analyse de graphe en flux, la
question n'est pas tant de savoir exécuter un algorithme que de savoir
\emph{ne pas} l'exécuter sur ce qui n'a pas changé.

La plateforme reste un prototype --- données synthétiques, un seul nœud,
latence de quelques secondes --- mais un prototype dont chaque limite est
mesurée et chaque piste d'amélioration identifiée. Elle montre qu'une
détection de structures de blanchiment explicable, auditable et proche du
temps réel est à la portée d'une pile entièrement open source.

% ============================================================
% WEBOGRAPHIE
% ============================================================
\begin{thebibliography}{99}
\phantomsection\addcontentsline{toc}{chapter}{Webographie}
\small
\bibitem{ibmaml} E.~Altman \emph{et al.}, \emph{Realistic Synthetic Financial Transactions for
Anti-Money Laundering Models}, NeurIPS 2023 (Datasets and Benchmarks).
\url{https://arxiv.org/abs/2306.16424}
\bibitem{anrf} Autorité Nationale du Renseignement Financier (Maroc). \url{https://www.anrf.gov.ma}
\bibitem{gafi} GAFI / FATF --- \emph{Jurisdictions under Increased Monitoring}.
\url{https://www.fatf-gafi.org}
\bibitem{spark} Apache Spark --- \emph{Structured Streaming Programming Guide} (3.5).
\url{https://spark.apache.org/docs/3.5.7/structured-streaming-programming-guide.html}
\bibitem{graphframes} GraphFrames --- graphes sur DataFrames Spark.
\url{https://github.com/graphframes/graphframes}
\bibitem{delta} Delta Lake --- documentation. \url{https://docs.delta.io}
\bibitem{deltars} delta-rs --- Delta Lake sans JVM (Rust, Python).
\url{https://delta-io.github.io/delta-rs/}
\bibitem{kafka} Apache Kafka --- documentation. \url{https://kafka.apache.org/documentation/}
\bibitem{streamlit} Streamlit --- documentation. \url{https://docs.streamlit.io}
\bibitem{pyvis} PyVis --- visualisation interactive de graphes. \url{https://pyvis.readthedocs.io}
\bibitem{pagerank} L.~Page, S.~Brin, R.~Motwani, T.~Winograd, \emph{The PageRank Citation Ranking:
Bringing Order to the Web}, Stanford InfoLab, 1999.
\bibitem{credits} Crédits graphiques : icônes Font Awesome Free 5 (licence CC BY 4.0,
\url{https://fontawesome.com}) ; logos des outils : Wikimedia Commons.
\end{thebibliography}
